% Table 1: Feature Ablation (event-likelihood loss, K=5, τ=0.1)
% Generated automatically — do not edit by hand
\begin{table}[t]
\centering
\caption{Feature ablation study under event-likelihood loss ($K=5,\tau=0.1$).
All metrics reported on the test split (mean over 3 seeds).
F3 removes ISAC BS-link residual statistics, retaining only GPS-ISAC position discrepancy
and signed neighbor anomaly contrast.}
\label{tab:feature_ablation}
\begin{tabular}{lcccccc}
\toprule
\multicolumn{7}{c}{\textbf{(a) RCS = 6 (practical conditions)}} \\
\midrule
Feature Set & Dim & Event F1$\uparrow$ & Event Prec.$\uparrow$ & Event Rec.$\uparrow$ & FAR$\downarrow$ & Delay (s)$\downarrow$ \\
\midrule
  Position only & 3 & 0.873 & 0.925 & 0.826 & 0.99 & 4.86 \\
  Local & 8 & 0.965 & 0.955 & 0.975 & 0.68 & 2.87 \\
  Full & 9 & 0.966 & 0.959 & 0.974 & 0.62 & 2.94 \\
\midrule
\multicolumn{7}{c}{\textbf{(b) RCS = 3 (degraded conditions)}} \\
\midrule
  Position only & 3 & 0.823 & 0.914 & 0.749 & 1.05 & 5.49 \\
  Local & 8 & 0.943 & 0.944 & 0.942 & 0.83 & 3.32 \\
  Full & 9 & 0.949 & 0.955 & 0.944 & 0.65 & 3.39 \\
\bottomrule
\end{tabular}
\end{table}


% Table 2: Step-level BCE vs. Event-likelihood Loss
% Generated automatically — do not edit by hand
\begin{table}[t]
\centering
\caption{Impact of event-likelihood loss compared to step-level BCE.
Both use the F9 feature set (ISAC + signed neighbor contrast).
Metrics on test split (mean over 3 seeds).}
\label{tab:loss_comparison}
\begin{tabular}{lcccccc}
\toprule
\multicolumn{7}{c}{\textbf{(a) RCS = 6 (practical conditions)}} \\
\midrule
Loss & Event F1$\uparrow$ & Event Prec.$\uparrow$ & Event Rec.$\uparrow$ & FAR$\downarrow$ & Delay (s)$\downarrow$ & Detect@2s$\uparrow$ \\
\midrule
  BCE (step-level) & 0.944 & 0.912 & 0.979 & 1.40 & 2.87 & 0.576 \\
  Event-likelihood & 0.966 & 0.959 & 0.974 & 0.62 & 2.94 & 0.574 \\
\midrule
\multicolumn{7}{c}{\textbf{(b) RCS = 3 (degraded conditions)}} \\
\midrule
  BCE (step-level) & 0.918 & 0.886 & 0.952 & 1.81 & 3.35 & 0.496 \\
  Event-likelihood & 0.949 & 0.955 & 0.944 & 0.65 & 3.39 & 0.493 \\
\bottomrule
\end{tabular}
\end{table}


% Table 3: Cross-Domain Generalization (event-likelihood loss, K=5)
% Generated automatically — do not edit by hand
\begin{table}[t]
\centering
\caption{Cross-domain generalization. Models trained on S-curve trajectories,
evaluated on multi-trajectory test sets with matched or mismatched RCS.
All metrics on test split (mean over 3 seeds).}
\label{tab:cross_domain}
\begin{tabular}{llcccc}
\toprule
Train Domain & Test Domain & Feat. & Event F1$\uparrow$ & FAR$\downarrow$ & Delay (s)$\downarrow$ \\
\midrule
\multicolumn{6}{c}{\textbf{Panel A: RCS=6 S-curve → RCS=6 Multi-traj (trajectory generalization)}} \\
\midrule
  RCS=6 S-curve & RCS=6 S-curve & F8 & 0.965 & 0.68 & 2.87 \\
  RCS=6 S-curve & RCS=6 S-curve & F9 & 0.966 & 0.62 & 2.94 \\
  RCS=6 S-curve & RCS=6 S-curve & F3 & 0.873 & 0.99 & 4.86 \\
\cmidrule{2-6}
  RCS=6 S-curve & RCS=6 Multi-traj & F8 & 0.937 & 1.89 & 2.00 \\
  RCS=6 S-curve & RCS=6 Multi-traj & F9 & 0.887 & 3.69 & 1.93 \\
\midrule
\multicolumn{6}{c}{\textbf{Panel B: RCS=6 S-curve → RCS=3 Multi-traj (cross-RCS + cross-trajectory)}} \\
\midrule
  RCS=6 S-curve & RCS=3 Multi-traj & F8 & 0.698 & 13.90 & 2.34 \\
  RCS=6 S-curve & RCS=3 Multi-traj & F9 & 0.605 & 19.16 & 2.31 \\
\midrule
\multicolumn{6}{c}{\textbf{Panel C: RCS=3 S-curve → RCS=6 Multi-traj (reverse cross-RCS)}} \\
\midrule
  RCS=3 S-curve & RCS=3 S-curve & F8 & 0.943 & 0.83 & 3.32 \\
  RCS=3 S-curve & RCS=3 S-curve & F9 & 0.949 & 0.65 & 3.39 \\
  RCS=3 S-curve & RCS=3 S-curve & F3 & 0.823 & 1.05 & 5.49 \\
\cmidrule{2-6}
  RCS=3 S-curve & RCS=6 Multi-traj & F8 & 0.968 & 0.81 & 2.09 \\
  RCS=3 S-curve & RCS=6 Multi-traj & F9 & 0.967 & 0.80 & 2.09 \\
\bottomrule
\end{tabular}
\end{table}
